\subsubsection{Planteamiento del problema}

Actualmente, el centro de salud de San Marcos enfrenta deficiencias significativas en la organización y el control de los turnos médicos, derivadas principalmente de la dependencia de métodos de registro manuales y obsoletos. El uso de papel y procesos no estandarizados ralentiza la recepción de los usuarios, propicia errores humanos y dificulta el manejo seguro y ágil de la información personal y clínica de los pacientes.

Como resultado directo de esta gestión ineficiente, se evidencia diariamente la formación de largas filas y tiempos de espera excesivamente elevados. Esta situación genera un notable desorden en el flujo de atención, lo que se traduce en frustración, desgaste físico y malestar para los pacientes. Paralelamente, el personal administrativo y médico sufre de sobrecarga laboral y estrés continuo al intentar gestionar grandes volúmenes de personas sin las herramientas adecuadas.

Una de las debilidades más críticas del modelo actual es la falta de un sistema automatizado que permita el registro, la asignación y, sobre todo, la priorización de turnos. Al carecer de un mecanismo que clasifique de manera objetiva a los pacientes según su nivel de urgencia médica o condición de vulnerabilidad (como adultos mayores, mujeres embarazadas o personas con discapacidad), se vulnera la equidad en el servicio y se compromete la atención oportuna de quienes más lo necesitan.
